\subsection{Two smaller finite detectors for the reduced central obstruction}

Continue with the marked discrete quotient
$B_0\to\mathcal T$ and kernel $L$.
The finite group $\mathcal T$ has the two specified
normal prime-power cores
\[
 \mathcal C_2=O_3(\mathcal T)\cong C_3,
       \qquad \mathcal C_3=O_2(\mathcal T)\cong Q_8.                    \tag{B.S44}
\]
The subscript is the coefficient prime in what follows.
Let $V^{(p)}$ be the inverse image of $\mathcal C_p$ in $B_0$,
and retain $T_2=\ker(B_0\to S_3)$ from B.2.10.
Define the actual normal subgroups
\[
\begin{aligned}
 P_2={}&L\cap[V^{(2)},V^{(2)}](V^{(2)})^2
              \cap[T_2,T_2]T_2^2,\\
 P_3={}&L\cap[V^{(3)},V^{(3)}](V^{(3)})^3.
\end{aligned}                                                   \tag{B.S45}
\]
All terms are normal in $B_0$, so these are literal
intersections, with no choice of a group identification.
Write $\overline P_p$ for their closures in $B$.
Put
\[
 Q_p=B/\overline P_p,
       \qquad R_p=U/\overline P_p^{\,6}.                       \tag{B.S46}
\]
The second map uses the six actual coordinates of
$N=L_0^6\subseteq U$.  Equivalently, $R_p$ is the
induced source image for the target $Q_p$ under the
same six maps as B.2.13.  Let $K_p$ be its kernel
in $B$, and put
\[
 R_p^{\mathrm c}=B/\overline{[B,K_p]},\qquad
 M_p=K_p/\overline{[B,K_p]}.
\]
The pure subgroup of $R_p$ is the image of $J$.
Let $S_p$ be generated in $M_p$ by the commutators
of lifts of its pure $s$- and $r$-centralizers with
the named torsion generators, exactly as in (B.S33).

\begin{result}{B.2.22}{Proved}{Finite detection of both remaining elementary obstructions}
The specified finite groups have orders
\[
\begin{array}{c|c|c|c}
 p&L/P_p&|Q_p|&|R_p|\\ \hline
 2&C_2^{11}&2^{15}3^2&2^{80}3^6\\
 3&C_3^4&2^4 3^6&2^{14}3^{30}
\end{array}                                                     \tag{B.S47}
\]
Their actual source quotient maps induce injective maps
\[
 \frac{\overline M}{p\overline M}
       \lhook\joinrel\longrightarrow
                 \frac{M_p}{pM_p+S_p},
                         \qquad p=2,3.                        \tag{B.S48}
\]
Here additive notation is used for the finite central
kernels.  Thus the two right-hand sides are finite
detectors of all remaining elementary reduced classes;
no class in the left-hand side is omitted by these maps.
\end{result}

\begin{proof}
We first identify the kernels in (B.S47) through their
actual homology maps.  The groups $V^{(2)}$ and
$V^{(3)}$ are torsion-free subgroups of indices 48 and
18 in $B_0$.  The torsion-freeness follows because the
named elements of orders two and three retain their
orders in $\mathcal T/\mathcal C_p$.  They are free of ranks
nine and four.  Since $[V^{(2)}:L]=3$ and
$[V^{(3)}:L]=8$ are prime to their respective coefficient
primes, their homology maps give exact identifications
\[
 H_1(V^{(p)},\mathbb F_p)
                    \cong (V_p)_{\mathcal C_p},
       \qquad \dim (V_2)_{\mathcal C_2}=9,
       \quad \dim (V_3)_{\mathcal C_3}=4.                              \tag{B.S49}
\]
This is the finite-index homology sequence with
coprime averaging, so it also specifies the map from
$H_1(L,\mathbb F_p)$.

At the prime two there is one further quotient.
The map from $L$ to $H_1(T_2,\mathbb F_2)$ identifies
that quotient with $(V_2)_{T_2}$, of dimension two.
Indeed the latter dimension is (B.S42), the former
dimension is two since $T_2$ is free of rank two,
and the homology sequence is onto since
$T_2/L\cong\operatorname{SL}_2(\mathbb F_3)$ has
abelianization $C_3$.  Its two-dimensional module is
the augmentation module of the three cusps of $T_2$.
The subgroup $\mathcal C_2=O_3(\mathcal T)$ acts on those
cusps as a three-cycle, so this module has no
$\mathcal C_2$ invariants.  Coprime averaging separates it
from $(V_2)_{\mathcal C_2}$, on which $\mathcal C_2$ acts trivially.
Consequently the joint map in (B.S45) is onto their
direct sum.  It gives the exact quotient modules
\[
 A_2=L/P_2=(V_2)_{\mathcal C_2}\oplus(V_2)_{T_2},
                 \qquad A_3=L/P_3=(V_3)_{\mathcal C_3}.                \tag{B.S50}
\]
These are elementary abelian, of dimensions eleven
and four.  Since $J$ is the full inverse image of
$J_{\mathcal T}$ in $B^6$, their sixfold kernels in
$R_p\to E$ are exactly $A_p^6$.  This proves every
order in (B.S47).

We record the exact cohomology comparison for the
first homology modules.  In the actual finite group
$J_{\mathcal T}\subseteq\mathcal T^6$, the normal
subgroups
\[
 D_p=J_{\mathcal T}\cap \mathcal C_p^6
\]
have orders $27$ for $p=2$ and $2048$ for $p=3$.
Their first-coordinate maps are onto $\mathcal C_p$.
These are small finite-group equalities checked
directly in
\path{work/centralizer_code/check_relation_coprime_cores.g}
and its executed output: the groups are the literal
intersections in the actual six-coordinate group,
with their specified coordinate maps.
Each $D_p$ is normal and has order prime to $p$.

The module $V_p$ splits into its $\mathcal C_p$ invariant
summand and a complementary summand with no $\mathcal C_p$
invariants, by the averaging idempotent.
Any module in the latter summand has zero cohomology
in every degree over $J_{\mathcal T}$: first take
$D_p$ cohomology, which is its invariant module in
degree zero and is zero in positive degrees, and
then use the quotient spectral sequence.  The
first-coordinate surjection makes that invariant
module zero.  At $p=2$ the additional two-dimensional
summand in (B.S50) also has this vanishing property.
It follows for the exact dual inclusions that
\[
 H^n(J_{\mathcal T},A_p^*)
             \xrightarrow{\ \sim\ }
 H^n(J_{\mathcal T},V_p^*)
                \quad\hbox{for every }n\geq0.                 \tag{B.S51}
\]
Duals are over $\mathbb F_p$.  This is an all-degree
comparison supplied by the actual normal coprime
subgroup, rather than an extrapolation from a
computed first-cohomology dimension.

Apply Shapiro to the six transitive factors as in
(B.S40).  Equation (B.S51) gives isomorphisms in
every column of the first-cohomology row of the
two Lyndon--Hochschild--Serre spectral sequences
for
\[
 A_p^6\longrightarrow R_p\longrightarrow E,
       \qquad L^6\longrightarrow U_0\longrightarrow E.
                                                                  \tag{B.S52}
\]
The maps on the degree-zero row are identities.
The map on the invariant part of the degree-two
row is onto.  At $p=3$ the target of that map is
zero by (B.S42).  At $p=2$ every surviving invariant
pair form factors through $(V_2)_{T_2}$ or
$(V_2)_{T_3}$ in the corresponding two coordinates.
The former is included in $A_2$; the latter is a
quotient of $(V_2)_{\mathcal C_2}$ because
$\mathcal C_2\subseteq T_3/L$.  Each such form therefore
has a unique lift as a cross-factor form on $A_2^6$.
It is invariant under $E$, since the quotient map
is equivariant and onto.  The cross-factor Kunneth
summands are preserved by the actual monomial
action, with their usual signed swaps, so they
define invariant second-cohomology classes on
$A_2^6$.

These row comparisons imply the concrete surjection
\[
 H^2(R_p,\mathbb F_p)
                   \twoheadrightarrow H^2(U_0,\mathbb F_p).
                                                                  \tag{B.S53}
\]
Here is the differential check.  The term in bidegree
$(2,0)$ has the same incoming differential from
$(0,1)$ in both sequences.  The term in bidegree
$(1,1)$ has the same outgoing differential to
$(3,0)$.  Thus their limiting terms map
isomorphically.  An invariant class in $(0,2)$
which survives on the right has an invariant lift
on the left.  Its second differential vanishes
because the maps on $(2,1)$ are isomorphisms.
Its third differential vanishes because the maps
on $(3,0)$ after quotienting by the second
differential from $(1,1)$ are isomorphisms.
There are no later differentials from $(0,2)$.
It therefore has a surviving lift.  The finite
filtration of total degree two proves (B.S53).

The abelianizations of $B_0$, $U_0$ and $R_p$ are
the same $C_6$.  The universal-coefficient sequence
and (B.S53) consequently give an injection from
$H_2(U_0,\mathbb Z)/p$ into the elementary quotient
of the canonical central kernel over $R_p$.
Equivalently, every marked central $C_p$ extension
of $U_0$ is the pullback of a marked central
$C_p$ extension of $R_p$; its restriction to the
two source torsion factors is retained by the
same map from $B_0$.

The actual profinite $M$ is a quotient of
$H_2(U_0,\mathbb Z)$ by B.2.15, and its map to
$M_p$ factors this last injection.  The resulting
map
\[
                         M/pM\lhook\joinrel\longrightarrow M_p/pM_p
                                                                  \tag{B.S54}
\]
is therefore injective as well: the preceding
surjection onto $M/pM$ followed by this map is
already injective.

Finally the pure centralizers of $s$ and $r$ in
$U$ map onto the corresponding centralizers in
$R_p$.  Their kernel is the product
$\overline P_p^{\,6}$, on which $s$ and $r$ again
have three two-cycles and two three-cycles without
fixed factors.  The same cycle-by-cycle cocycle
equation as in B.2.20 lifts every centralizing
element.  Hence the image of $S$ in $M_p$ is
exactly $S_p$.  Quotienting the injection (B.S54)
by these actual image subgroups gives precisely
(B.S48).
\end{proof}
